% archon:covers Gluck/Sphere.lean
% archon:covers Gluck/Sphere/Defs.lean Gluck/Sphere/Flow.lean Gluck/Sphere/Admissible.lean
% archon:covers Gluck/Sphere/ArcAlgebra.lean Gluck/Sphere/StepReparam.lean Gluck/Sphere/Margins.lean
% archon:covers Gluck/Sphere/FirstVariation.lean Gluck/Sphere/Reconstruction.lean
% archon:covers Gluck/Sphere/ConjWinding.lean Gluck/Sphere/EndpointWinding.lean Gluck/Sphere/Converse.lean
\chapter{Spherical geometry and definitions}
\label{chap:sphere}

\paragraph{Formalization map.}
This chapter corresponds primarily to \texttt{Gluck/Sphere/Defs.lean}; the
subsequent spherical chapters cover the flow, winding and converse modules.

This chapter formalizes the \emph{spherical} converse to the four vertex theorem in
the positive-curvature (Gluck-style) regime: a prescribed cyclic geodesic-curvature
function on the round sphere $S^2$, with the four-vertex property and a uniform lower
bound forcing confinement to an open hemisphere, is the geodesic curvature of a simple
closed curve contained in that hemisphere. The Euclidean development
(\cref{chap:reduction} and its dependencies, culminating in \texttt{gluck\_converse} and
\texttt{dahlbergConverse}) is complete and is reused, not disturbed. The positive-stage
capstone is named \texttt{sphericalConverse\_pos}; the name \texttt{sphericalConverse} is
reserved for the mandatory mixed-sign stage 2.

All spherical geometry is transported to the \emph{stereographic disk model}: the open
unit disk $\{|z|<1\}\subset\mathbb C$ with the round metric
$g_{S^2}=\dfrac{4}{(1+|z|^2)^2}\,|dz|^2$. Stereographic projection from the missed
antipode is a conformal diffeomorphism onto $S^2$ minus a point, so a curve confined to
$\{|z|<1\}$ is simple on $S^2$ iff its disk representative is simple in $\mathbb C$.

% SOURCE: references/spaceform_notes.md, "THE CRUX AND ITS RESOLUTION" +
%   "H² vs S²" sections (read from references/spaceform_notes.md)
% SOURCE QUOTE: "The tangent-angle gauge reduces closure to an honest 2-D problem.
%   Parametrize by the EUCLIDEAN tangent angle θ in the disk model, z_θ = q(θ)(cos θ,
%   sin θ), q>0. Then κ_E = 1/q and the conformal law solves ALGEBRAICALLY for the
%   speed:  S²:  q(θ) = (1+|z|²) / [2(κ_S(θ) + ⟨z,n⟩)]  (den. >0 ⟺ κ_S > −⟨z,n⟩).
%   Reconstruction becomes a nonlinear 2-D ODE z_θ = q(θ)(cos θ, sin θ) on the disk.
%   The tangent angle increases once around ⇒ TURNING NUMBER +1 IS BUILT IN ⇒ the
%   rotational/framing dimension is discharged by construction ..., leaving only the
%   2-D point-closure error E(z₀)=z(2π;z₀)−z₀ ∈ ℝ². Planar winding/degree then applies."

\section{The conformal geodesic-curvature law}
\label{sec:sphere_conformal}

Write $z=(x,y)\in\mathbb R^2\cong\mathbb C$ for a point of the disk and let a curve have
unit Euclidean tangent $T=e^{i\varphi}$ at inclination angle $\varphi$. We use the
\emph{outward} (right) unit normal
\[
  n=-iT=e^{i(\varphi-\pi/2)}=(\sin\varphi,-\cos\varphi),\qquad
  \langle z,n\rangle=x\sin\varphi-y\cos\varphi .
\]
For the isothermal factor $e^{2\lambda}=4/(1+|z|^2)^2$ one has
$\lambda=\ln 2-\ln(1+|z|^2)$ and $\nabla\lambda=-2z/(1+|z|^2)$, so
$\partial_n\lambda=\langle\nabla\lambda,n\rangle=-2\langle z,n\rangle/(1+|z|^2)$. The
conformal transformation law $\kappa_g=e^{-\lambda}(\kappa_E+\partial_n\lambda)$ gives the
fundamental identity used throughout:
\[
  \kappa_S \;=\; \frac{1+|z|^2}{2}\,\kappa_E \;-\; \langle z,n\rangle .
\]
Here $\kappa_E$ is the Euclidean signed curvature of the disk representative and
$\kappa_S$ its spherical geodesic curvature. In the tangent-angle gauge (\cref{sec:sphere_defs})
$\kappa_E=\varphi'/\|z'\|$, so the law rearranges to the \emph{speed relation}
\[
  \frac{1+|z|^2}{2}\,\varphi' \;=\; \bigl(\kappa_S+\langle z,n\rangle\bigr)\,\|z'\| ,
  \qquad\text{gauge speed}\quad q=\frac{1+|z|^2}{2\bigl(\kappa_S+\langle z,n\rangle\bigr)} .
\]

\paragraph{Geodesic-circle sanity anchor (decisive sign check).}
A counterclockwise Euclidean circle $|z|=r$ has $\kappa_E=1/r$ and outward normal
$n=z/r$, so $\langle z,n\rangle=r$ and
$\kappa_S=\dfrac{1+r^2}{2r}-r=\dfrac{1-r^2}{2r}=\cot\rho$ with $r=\tan(\rho/2)$ — exactly
the geodesic curvature of a spherical circle of intrinsic radius $\rho$ (and $0$ on the
equator $\rho=\pi/2$). The gauge speed is then $q=r$. The opposite sign choice would give
$(1+3r^2)/(2r)\ne\cot\rho$, so \emph{the outward normal with a minus sign is the correct
convention}; any Lean formalization must reproduce $q(\text{circle }r)=r$.

\gluckfigure{sphere-stereographic}
  {The spherical conformal-curvature law in the stereographic disk. The centered
   circle of radius $r=1/2$ realizes the sign-check value
   $\kappa_S=(1-r^2)/(2r)=3/4$ and has gauge speed $q=r$.}
  {A circle inside the stereographic unit disk with a marked point, tangent
   arrow, and outward normal. Beside it are the conformal metric, curvature law,
   and the evaluated curvature three quarters.}
  {fig:spherical_stereographic_sign_check}

\paragraph{Lean encoding of the normal term.}
In Lean the left normal $iT$ is the natural object: $\langle z,iT\rangle=\langle z,
i\,e^{i\varphi}\rangle=-x\sin\varphi+y\cos\varphi=-\langle z,n\rangle$. Hence the speed
relation is encoded with $\langle z,n\rangle$ replaced by $-\langle z,iT\rangle$:
\[
  \frac{1+\|z\|^2}{2}\,\varphi'
    =\bigl(\kappa-\langle z,\, i\,e^{i\varphi}\rangle_{\mathbb R}\bigr)\,\|z'\| .
\]
The project supplies the rewrite lemma
$\langle z, i\,e^{i\varphi}\rangle_{\mathbb R}=-(\operatorname{Re}z)\sin\varphi
+(\operatorname{Im}z)\cos\varphi$ so both forms are interchangeable.

\section{Definition layer}
\label{sec:sphere_defs}

\begin{lemma}\leanok
[normal-term coordinate identity]
  \label{lem:sphere_normal_inner_eq}
  \lean{Gluck.sphereNormal_inner_eq}
  For $z\in\mathbb C$ and $\varphi\in\mathbb R$,
  $\langle z,\, i\,e^{i\varphi}\rangle_{\mathbb R}
   =-(\operatorname{Re}z)\sin\varphi+(\operatorname{Im}z)\cos\varphi$.
  Equivalently $\langle z, i\,e^{i\varphi}\rangle_{\mathbb R}=-\langle z,n\rangle$ for the
  outward normal $n=-i\,e^{i\varphi}$, so the defining coefficient of
  \cref{def:realizes_spherical_curvature} is
  $\kappa-\langle z,i\,e^{i\varphi}\rangle_{\mathbb R}=\kappa+\langle z,n\rangle$.
\end{lemma}

\begin{proof}

\leanok
  Expand the real inner product $\langle w,v\rangle_{\mathbb R}=(v\,\overline w).\mathrm{re}$
  with $v=i\,e^{i\varphi}$, using
  $(e^{i\varphi}).\mathrm{re}=\cos\varphi$, $(e^{i\varphi}).\mathrm{im}=\sin\varphi$; the
  real and imaginary multiplications collapse to
  $-(\operatorname{Re}z)\sin\varphi+(\operatorname{Im}z)\cos\varphi$.
\end{proof}

\begin{definition}\leanok
[realizes spherical curvature]
  \label{def:realizes_spherical_curvature}
  \lean{Gluck.RealizesSphericalCurvature}
  A curve $z:\mathbb R\to\mathbb C$ \emph{realizes the spherical curvature function}
  $\kappa:\mathbb R\to\mathbb R$ when it is $C^1$, regular ($z'(t)\neq 0$ for all $t$),
  confined to the open disk ($\|z(t)\|<1$ for all $t$), and there is a differentiable
  tangent-angle function $\varphi:\mathbb R\to\mathbb R$ with, for all $t$,
  \[
    z'(t)=\|z'(t)\|\,e^{i\varphi(t)},\qquad
    \frac{1+\|z(t)\|^2}{2}\,\varphi'(t)
      =\bigl(\kappa(t)-\langle z(t),\,i\,e^{i\varphi(t)}\rangle_{\mathbb R}\bigr)\,\|z'(t)\| .
  \]
  By \cref{sec:sphere_conformal} the coefficient $\kappa-\langle z,iT\rangle
  =\kappa+\langle z,n\rangle$ is the reciprocal-speed weight of the conformal law, so this
  says precisely that $\kappa$ is the spherical geodesic curvature of $z$. It mirrors the
  Euclidean \cref{def:realizes_curvature} exactly: it is \emph{speed-free} (only the angle
  $\varphi$, not $z$, must be differentiable; the speed $\|z'\|$ is never normalized to arc
  length), so it is meaningful for a merely continuous $\kappa$ and does not reintroduce
  the three-dimensional spherical frame problem.
\end{definition}

\begin{definition}\leanok
[spherical four-vertex condition, positive stage]
  \label{def:sphere_four_vertex}
  \lean{Gluck.SphereFourVertex}
  \uses{def:four_vertex_condition, def:is_curvature_function}
  A function $\kappa:\mathbb R\to\mathbb R$ satisfies the \emph{positive-stage spherical
  four-vertex condition} when it is a curvature function (\cref{def:is_curvature_function}:
  continuous, $2\pi$-periodic, strictly positive) and satisfies the value-separated Euclidean
  four-vertex condition \cref{def:four_vertex_condition} — exactly the hypothesis package of
  the Euclidean \texttt{gluck\_converse}. No confinement bound is bundled: by compactness a
  uniform bound $R$ with $0<R<1$ and $\kappa>R$ is automatic
  (\cref{lem:sphere_curvature_lower_bound}), and a bundled witness would be gratuitous
  interface drift from the Euclidean theorem without carrying any datum the realization
  layer consumes (audit finding, iter-043). The bound $\kappa>R$ is what makes the
  denominator $\kappa_S+\langle z,n\rangle$ strictly positive on $\{|z|\le R\}$
  (\cref{lem:sphere_denom_pos}); it is the positive-stage stand-in for the sign-changing
  admissibility $\kappa_S>-\langle z,n\rangle$ of stage 2.
\end{definition}

\begin{lemma}\leanok
[uniform curvature lower bound]
  \label{lem:sphere_curvature_lower_bound}
  \lean{Gluck.exists_curvature_lower_bound}
  \uses{def:is_curvature_function}
  If $\kappa$ is a curvature function (\cref{def:is_curvature_function}), then there is $R$
  with $0<R$, $R<1$, and $R<\kappa(\theta)$ for all $\theta$.
\end{lemma}

\begin{proof}

\leanok
  $\kappa$ is continuous on the compact interval $[0,2\pi]$, so it attains there a minimum
  value $m=\kappa(\theta_0)>0$. Set $R=\min(m,1)/2$. Then $0<R$, $R<1$, and
  $R<m\le\kappa(\theta)$ for every $\theta\in[0,2\pi]$; an arbitrary $\theta$ reduces to
  this window by $2\pi$-periodicity (subtract the integer multiple
  $2\pi\lfloor\theta/(2\pi)\rfloor$).
\end{proof}

\begin{theorem}\leanok
[spherical converse, positive stage]
  \label{thm:spherical_converse_pos}
  \lean{Gluck.sphericalConverse_pos}
  \uses{def:realizes_spherical_curvature, def:sphere_four_vertex,
        lem:spherical_speed_solve, lem:reconstruction_ode,
        lem:spherical_endpoint_winding, lem:realizes_spherical_comp,
        lem:spherical_simplicity, lem:spherical_circle_realizes}
  If $\kappa$ satisfies the positive-stage spherical four-vertex condition
  (\cref{def:sphere_four_vertex}), then there is a simple closed curve $z$ confined to the
  open disk that realizes $\kappa$ as its spherical geodesic curvature:
  \[
    \texttt{SphereFourVertex }\kappa \;\Longrightarrow\;
    \exists\,z,\ \texttt{IsSimpleClosed }z\ \wedge\ \texttt{RealizesSphericalCurvature }z\,\kappa .
  \]
  This is the same conclusion shape as the Euclidean \texttt{gluck\_converse}, with
  \texttt{RealizesCurvature} replaced by its spherical analogue.
\end{theorem}

\begin{proof}
  \uses{lem:reconstruction_ode, lem:spherical_endpoint_winding,
        lem:spherical_simplicity, lem:spherical_circle_realizes,
        lem:realizes_spherical_comp, lem:exists_c1_circle_inverse,
        lem:is_simple_closed_comp}
        \leanok
  Split on the four-vertex disjunction (\cref{def:four_vertex_condition}). The
  \emph{constant} branch is discharged directly by the explicit centered circle
  (\cref{lem:spherical_circle_realizes}), with no flow machinery. For the value-separated
  (non-constant) branch: \cref{lem:spherical_endpoint_winding} supplies $R<1$,
  $\delta>0$, a circle reparametrization $h_1$, and a \emph{closed, admissible}
  trajectory $z$ of the $(\kappa\circ h_1,R,\delta)$-truncated flow.
  \cref{lem:reconstruction_ode} (with curvature function $\kappa\circ h_1$) upgrades its
  periodic extension to a closed $C^1$ curve confined to $\{|z|\le R\}\subset\{|z|<1\}$
  realizing $\kappa\circ h_1$ in the gauge $\varphi(\theta)=\theta$.
  \cref{lem:realizes_spherical_comp} with $\psi=h_1^{-1}$ pulls this back to a curve
  realizing $\kappa$ (closedness and injectivity-on-a-period transfer along the
  circle-equivariant $h_1^{-1}$, as in the Euclidean
  \texttt{isSimpleClosed\_comp}). \cref{lem:spherical_simplicity} gives simplicity in
  the disk, hence on $S^2$. This mirrors how \texttt{dahlbergConverse} assembles its
  closing parameter, simplicity transport, and reconstruction. \emph{(Assembled formally
  in S2-E once the S2-D layer exists in Lean.)}
\end{proof}

\section{Proof arc}
\label{sec:sphere_arc}

The proof follows the Euclidean Route A architecture, with the reconstruction integral
replaced by a reconstruction ODE. The architecture was pinned at iteration 046
(design analysis + external consult, load-bearing computations re-verified in-project):
\begin{enumerate}
\item \textbf{S2-B (truncated flow, unconditional).} The gauge speed is truncated
  algebraically (\cref{def:truncated_speed}) so the ODE field becomes globally
  Lipschitz and bounded; global existence on $[0,2\pi]$, uniqueness, and continuous
  dependence on the initial point are then unconditional, with no confinement
  prerequisite. Only the closed curve found at the end must be shown, a posteriori, to
  be \emph{admissible} (both truncations inactive), whereupon it is an honest spherical
  curve.
\item \textbf{S2-C (admissibility by margin transport).} The confinement mechanism is
  \emph{perturbative}: an explicit model trajectory (a spherical bicircle made of
  circular arcs, \cref{lem:constant_curvature_arc}) is admissible with quantitative
  margins, and a Gr\"onwall estimate with $L^1$-in-$\theta$ drive
  (\cref{lem:invariant_admissible_domain}) transports the margins to every trajectory
  whose curvature is $L^1$-close and whose start is near the model start. The
  Dahlberg-style reparametrization already formalized in \cref{chap:dahlberg_step1}
  supplies arbitrarily small $L^1$ distance to a two-value step model, so no
  uniform-in-$\kappa_0$ direct confinement estimate is needed (none exists: no disk is
  flow-invariant, see the discussion in \cref{lem:invariant_admissible_domain}).
\item \textbf{S2-D (endpoint map and winding).} The degree argument runs over a compact
  $2$-disk of \emph{initial points} $z_0$, but its computable model is the symmetric
  two-value ($\pi$-periodic) step curvature, exploited through the half-turn
  equivariance $z\mapsto -z$, $\theta\mapsto\theta+\pi$
  (\cref{lem:flow_half_turn_equivariance}). The constant-curvature model is
  \emph{unusable}: for constant $\kappa$ every admissible trajectory closes (all
  spherical circles), so the endpoint error vanishes identically and carries no degree
  — see the DO-NOT recorded in \cref{lem:spherical_endpoint_winding}.
\end{enumerate}
The speed, flow, transport, and step-model layers (S2-A through S2-D except the final
assembly) are formalized axiom-clean; what remains is the assembly
\cref{lem:spherical_endpoint_winding}, the S2-E simplicity chain, and the capstone.

\begin{definition}\leanok
[gauge speed]
  \label{def:spherical_speed}
  \lean{Gluck.sphericalSpeed}
  \uses{lem:sphere_normal_inner_eq}
  For $\kappa:\mathbb R\to\mathbb R$, $\theta\in\mathbb R$ and $z\in\mathbb C$ set
  \[
    q_\kappa(\theta,z)\;=\;\frac{1+\|z\|^2}
      {2\bigl(\kappa(\theta)-\langle z,\,i\,e^{i\theta}\rangle_{\mathbb R}\bigr)} .
  \]
  By \cref{lem:sphere_normal_inner_eq} the bracket equals
  $\kappa(\theta)+\langle z,n\rangle$ for the outward normal $n=-i\,e^{i\theta}$ of the
  gauge $\varphi(\theta)=\theta$, so $q_\kappa$ is the algebraic solution of the speed
  relation of \cref{def:realizes_spherical_curvature} for the speed $\|z'\|$. It is a total
  function of the junk-value kind: where the bracket vanishes it returns
  division-by-zero junk, and every lemma about it carries the admissibility hypothesis
  $\|z\|\le R<\kappa(\theta)$.
\end{definition}

\begin{lemma}\leanok
[denominator positivity]
  \label{lem:sphere_denom_pos}
  \lean{Gluck.sphere_denom_pos}
  \uses{lem:sphere_normal_inner_eq}
  If $\|z\|\le R$ and $R<\kappa(\theta)$, then
  $\kappa(\theta)-\langle z,\,i\,e^{i\theta}\rangle_{\mathbb R}\;\ge\;\kappa(\theta)-R\;>\;0$.
\end{lemma}

\begin{proof}

\leanok
  Cauchy--Schwarz gives
  $\bigl|\langle z,\,i\,e^{i\theta}\rangle_{\mathbb R}\bigr|\le\|z\|\cdot\|i\,e^{i\theta}\|
  =\|z\|\le R$, since $\|i\,e^{i\theta}\|=|i|\cdot|e^{i\theta}|=1$. Hence
  $\kappa(\theta)-\langle z, i\,e^{i\theta}\rangle_{\mathbb R}\ge\kappa(\theta)-R>0$.
\end{proof}

\begin{lemma}\leanok
[positive speed / algebraic solve]
  \label{lem:spherical_speed_solve}
  \lean{Gluck.sphericalSpeed_pos}
  \uses{def:spherical_speed, lem:sphere_denom_pos, lem:sphere_curvature_lower_bound}
  If $\|z\|\le R$ and $R<\kappa(\theta)$, then $q_\kappa(\theta,z)>0$ (the hypotheses are
  supplied under \cref{def:sphere_four_vertex} by \cref{lem:sphere_curvature_lower_bound}).
  Positivity of $q_\kappa$ makes the tangent turn once around as $\theta$ runs over
  $[0,2\pi]$, so along any single closed reconstructed curve the turning number is $+1$ by
  construction and the rotational/framing dimension is discharged (the exact role
  $\oint\kappa=2\pi$ plays in the plane). This fixes the framing dimension only; it does not
  by itself give point closure, confinement, or the boundary winding of the endpoint map.
\end{lemma}

\begin{proof}

\leanok
  The numerator $1+\|z\|^2$ is strictly positive, and the denominator
  $2(\kappa(\theta)-\langle z, i\,e^{i\theta}\rangle_{\mathbb R})$ is strictly positive by
  \cref{lem:sphere_denom_pos}; a quotient of positives is positive.
\end{proof}

\begin{lemma}\leanok
[speed continuity]
  \label{lem:spherical_speed_continuous}
  \lean{Gluck.sphericalSpeed_continuousOn}
  \uses{def:spherical_speed, lem:sphere_denom_pos}
  If $\kappa$ is continuous and $R<\kappa(\theta)$ for all $\theta$, then
  $(\theta,z)\mapsto q_\kappa(\theta,z)$ is continuous on the slab
  $\{(\theta,z):\|z\|\le R\}\subseteq\mathbb R\times\mathbb C$.
\end{lemma}

\begin{proof}

\leanok
  The numerator $(\theta,z)\mapsto 1+\|z\|^2$ is continuous, and the denominator
  $(\theta,z)\mapsto 2(\kappa(\theta)-\langle z, i\,e^{i\theta}\rangle_{\mathbb R})$ is
  continuous ($\kappa$ is continuous; the inner product is continuous bilinear and
  $\theta\mapsto i\,e^{i\theta}$ is continuous). On the slab the denominator is nonvanishing
  by \cref{lem:sphere_denom_pos}, so the quotient is continuous there.
\end{proof}

\begin{lemma}\leanok
[speed periodicity]
  \label{lem:spherical_speed_periodic}
  \lean{Gluck.sphericalSpeed_periodic}
  \uses{def:spherical_speed}
  If $\kappa$ is $2\pi$-periodic, then $q_\kappa(\theta+2\pi,z)=q_\kappa(\theta,z)$ for all
  $\theta$ and $z$.
\end{lemma}

\begin{proof}

\leanok
  $\kappa(\theta+2\pi)=\kappa(\theta)$ by hypothesis, and
  $e^{i(\theta+2\pi)}=e^{i\theta}e^{2\pi i}=e^{i\theta}$ since $e^{2\pi i}=1$; both
  numerator and denominator of $q_\kappa$ are unchanged.
\end{proof}

\begin{lemma}

\leanok
[geodesic-circle sanity anchor]
  \label{lem:spherical_speed_circle}
  \lean{Gluck.sphericalSpeed_circle}
  \uses{def:spherical_speed}
  For $r>0$ let $\kappa\equiv(1-r^2)/(2r)$ be the constant geodesic curvature of the
  spherical circle of stereographic radius $r$, and let $z=-r\,i\,e^{i\theta}$ be the point
  of that circle at tangent angle $\theta$ (counterclockwise gauge). Then
  $q_\kappa(\theta,z)=r$: the gauge speed reproduces the Euclidean radius. This is the
  decisive machine-checked \emph{sign certificate} of \cref{sec:sphere_conformal} — with
  the opposite sign convention the value would be $\ne r$ and the build would fail.
\end{lemma}

\begin{proof}

\leanok
  $\|z\|=r$ and $\langle z,\,i e^{i\theta}\rangle=-r\,\|i e^{i\theta}\|^2=-r$, so the
  denominator is $2(\kappa+r)=2\bigl(\tfrac{1-r^2}{2r}+r\bigr)=\tfrac{1+r^2}{r}$ and
  $q=\dfrac{1+r^2}{(1+r^2)/r}=r$.
\end{proof}
